En esta sección se detallan los entregables correspondientes al análisis de los dominios mediante el script en bash. Los dominios a analizar son: \texttt{unam.mx}, \texttt{pemex.com}, \texttt{gob.mx}, \texttt{ipn.mx} y una página a elección.

\subsection*{Requisitos}
% Se debe incluir todo lo que se necesita para obtener la información mediante la prueba de penetración.
% Por ejemplo ¿qué problema se quiere resolver o qué se quiere saber? ¿Qué información es necesaria para ello? ¿Para qué?
Dado que buscamos encontrar vulnerabilidades en los sistemas, requerimos recopilar la mayor cantidad de información disponible que nos proporcione una visión general de la infraestructura del sistema. De esta manera, en la práctica obtenemos información de conectividad y rendimiento, información legal y de registro, mapeo de direcciones IP y segmentos de red, la topología de tráfico y enrutamiento, la enumeración completa de DNS y el mapeo de subdominios, así como la información del dominio asociado al certificado SSL/TLS y las cabeceras HTTP de redirección en el caso de que usemos IPs para poder hacer el análisis completo con el script utilizando IPs como entrada, pues algunas herramientas son sensibles a utilizar únicamente nombres de dominio. Con el análisis de esta información podemos llegar a encontrar puntos débiles de los dominios o IPs para que, a partir de una búsqueda exhaustiva de vulnerabilidades en bases de datos oficiales, podamos en algún momento informar a las autoridades correspondientes (la información para estos reportes también se incluye en la búsqueda del script) para que se puedan implementar soluciones.
